\section{Existing Systems, Products and Techniques}\label{sec:existing-systems}

This section examines the systems, products and techniques that currently occupy
the space in which the Intent Recognition Engine operates. They are reviewed
against the same inclusion criteria as the wider survey (Section~\ref{sec:review-method-and-scope}): relevance to at
least one research question and a verifiable primary source reporting either an
empirical evaluation or an authoritative specification. Where a surveyed option
is a technique rather than a delivered product, the assessment asks whether it
can be combined with the deployment constraint adopted in this
dissertation---offline, CPU-only, single-inference, retrieval-free operation.

\subsection*{Deterministic security tooling}

The techniques that the IRE ultimately drives are mature and deterministic. Nmap
is the canonical network-discovery and scanning tool, whose techniques---host
discovery, port enumeration, service and version detection---define the target
values an intent contract must carry \parencite{lyon2009nmap}. Metasploit is the
canonical exploitation-framework driver, whose modules and options parameterise
the most privileged class of intents \parencite{metasploit}. Both expose
structured command interfaces (a tool name, target, options and constraints)
rather than free-form prose, and neither embeds a natural-language layer of its
own. In current practice the mapping from operator intent onto these structured
interfaces is performed by a human analyst or by a scripted wrapper, and any
authorisation policy has to be wired externally; RFC~1918 \parencite{rfc1918}
provides only the technical basis for distinguishing internal from external
addressing. These tools therefore constitute the deterministic ``ground truth''
that a translation layer must produce---and the reason a translation error
(Chapter~\ref{ch:product-design}) becomes a machine-dispatched error.

\subsection*{Language-model copilots for security}

The recent systems that add natural language to this toolchain are generation-side
copilots. PENTESTGPT decomposes a penetration-testing engagement across reasoning,
generation and guidance chains and executes each phase incrementally, producing
structured guidance that a human operator reviews \parencite{deng2024pentestgpt}.
The sparring-partner system demonstrated by Happe and Cito is a closed loop in
which a model proposes low-level actions that are executed automatically against
a deliberately vulnerable machine, up to root access \parencite{happe2023pwned}.
Both establish that a model can drive security tooling in controlled
environments, and both treat output correctness informally: the intended consumer
of the output is human, and neither system subjects generation to deterministic
schema, scope or authorisation validation before an action is taken. As products,
they are representative of the ``AI assistant'' category of security copilots
rather than of an automated, machine-consumed intent layer.

\subsection*{Model-integration techniques: retrieval and tool calling}

Beyond security copilots, two integration techniques are common to LLM-driven
automation generally. Retrieval-augmented generation (RAG) conditions answers on
external corpora so that claims can be traced to sources
\parencite{lewis2020rag}, using sentence embeddings as the similarity signal
\parencite{reimers2019sbert}; tool-calling scaffolds connect a model to external
commands through function-call interfaces. As deployed techniques both widen the
integration surface: retrieved content enters the prompt at inference time and is
therefore an indirect-injection vector \parencite{greshake2023indirect}, and
tool-calling binds a model action to a machine effect with no intermediate
validation stage. The IRE is deliberately retrieval-free and does not render
free-form model output to a downstream shell, which narrows these channels to the
classification of a single utterance (Section~\ref{sec:domain-foundations}).

\subsection*{Guardrail techniques for model-integrated applications}

At the application layer, guardrail practice is codified by the OWASP Top~10 for
LLM applications, which catalogues prompt injection (LLM01), insecure output
handling and related classes as first-class risks rather than theoretical
concerns \parencite{owasp2025}. The defence techniques surveyed in
Section~\ref{sec:related-research} are the prevention-detection families
benchmarked by Liu \emph{et al.} \parencite{liu2023promptinjection}---shown in
that study to be individually insufficient---together with detection techniques
for indirect injection \parencite{greshake2023indirect} and the empirically
demonstrated bypass of deployed content safeguards by dual-use attacks
\parencite{kang2024dualuse}. What these techniques have in common is that they
act as app-layer checks over a model's behaviour after the fact. None of them is
integrated as a deterministic validation stack that sits \emph{between} model
generation and machine dispatch and that makes syntax, scope and authorisation
independently decidable (Section~\ref{sec:domain-foundations}).

\subsection*{Fabrication-detection and uncertainty techniques}

A separate product-facing technique family detects fabricated output. SelfCheckGPT
scores the mutual consistency of sampled responses \parencite{manakul2023selfcheckgpt};
semantic-uncertainty and semantic-entropy methods measure answer disagreement over
meaning-equivalence clusters \parencite{kuhn2023semantic,farquhar2024semantic};
and self-consistency decoding accumulates agreement across samples
\parencite{wang2023selfconsistency}. These are deployed as techniques that assume
repeated or interactive generation. Two external references complement them for
the value spaces relevant here: the National Vulnerability Database (NVD) is the
authoritative registry of CVE identifiers \parencite{nistnvd}, and the Exploit
Prediction Scoring System (EPSS) scores the likelihood that a vulnerability will
be exploited \parencite{firstepss}. Both are the deterministic, offline-unfriendly
references that a fabricated-identifier guard would consult if the offline
boundary were relaxed (Chapter~\ref{ch:results}).

\subsection*{Comparative assessment}

Table~\ref{tab:existing-systems} consolidates the judgments from the subsections above against the properties that matter for the RQ1--RQ3
measurement (Section~\ref{sec:research-questions}). ``Schema check'' and
``scope/authorisation enforcement'' record whether a system subjects
machine-consumed output to structural and policy checks; ``offline CPU'' and
``retrieval-free'' record whether the system can run under the deployment
constraint adopted here.

\begin{table}[h]
\centering
\footnotesize
\caption{Surveyed existing systems and techniques against the properties adopted in this dissertation. Judgements summarise the cited primary sources and are re-verified at the final AUDIT.}
\label{tab:existing-systems}
\begin{tabular}{@{}lcccc@{}}
\hline
\textbf{System / technique} & \textbf{Human review} & \textbf{Schema check} & \textbf{Scope / authz} & \shortstack{\textbf{Offline CPU}\\\textbf{retrieval-free}} \\
\hline
Nmap / Metasploit & -- & Yes (interface) & External only & Yes \\
PENTESTGPT & Yes & No & No & -- \\
Sparring partner & Partial & No & No & -- \\
RAG agents & Partial & No & No & No \\
Fabrication detection & -- & No & No & Partial \\
IRE (this work) & No & Yes & Yes & Yes \\
\hline
\end{tabular}
\end{table}

In summary, the surveyed landscape contains tools that are deterministic but
require manual intent mapping, copilots that map intent but consume human review,
integration techniques that connect models to tools without an intermediate
validation stage, and defence techniques that check model behaviour after the
fact. None combines a deterministic schema, scope and authorisation validation
stack with offline, CPU-only, single-inference and retrieval-free operation over
a closed intent taxonomy. That combination is the space the existing-systems
survey leaves open and is consolidated in Section~\ref{sec:gap-synthesis}.

\TODO{If the assessment rubric requires a survey of \emph{named} commercial
security copilots (for example cloud SIEM/EDR copilot products), provide
verifiable primary sources for each; none are currently in the reference
database, so no commercial product is named here.}

\TODO{At the final AUDIT step, re-verify the comparative-table judgements in
Table~\ref{tab:existing-systems} against the full text of each cited source.}